\section{Introduction}\label{sec:intro}
Time series learning is a cornerstone in numerous modern data mining pipelines where related tasks involve forecasting, classification, anomaly detection, and decision-making in dynamical environments~\cite{esling2012timeseries, chandola2009anomaly, aaraba2024fr}. The main challenge is not only fitting nonlinear temporal dependencies, but also doing so under practical limitations such as limited labeled data, computational budgets, and the need for stable representations when the underlying process exhibits dependence and nonstationary effects. In this direction, reservoir computing~\cite{Jaeger2001EchoState, Grigoryeva_Ortega_2018} has been an attractive modeling approach as it delegates most of the representational burden to a rich dynamical system, while restricting training to a lightweight readout, which yields a favourable trade-off between expressivity and optimization cost.

Quantum reservoir computing (QRC) extends this paradigm by leveraging quantum dynamics to generate features for sequential inputs~\cite{Fujii_Nakajima_2017,Chen_Nurdin_2019,Chen_Nurdin_Yamamoto_2020}. %\cite{Suzuki_2022,Mujal_2023,Hu_2024}. 
By injecting temporal data into a quantum system and subsequently learning a classical readout, QRC proposes a principled way of representing time series with potentially rich, compact internal quantum states. However, there are two obstacles that limit the adoption of this approach as a reliable data mining tool. The first concerns the fact that existing designs are either inconvenient to implement faithfully on hardware or rely on expensive measurement schemes~\cite{mujal2023weakproj, zhu2025minimalistic}. The second is that works on the subject often lack end-to-end frameworks that connect concrete, efficiently implementable QRC architectures to explicit learning-theoretic guarantees on the generalization capacity of such methods~\cite{mujal2021opportunities, qrc_risk_bounds_2025}.

To reduce this gap, we propose in this work an implementable QRC architecture called \textsc{QuaRK} that is systematically designed with scalability in mind. First, a Johnson--Lindenstrauss (JL)~\cite{mohri2018foundations} projection decouples the quantum resources (quantum circuit width and depth) from the original data dimension, enabling the model to remain practical even in the case of high-dimensional inputs. Each time series point is projected to this dimension, yielding a different representation of the initial series with dimensionality matching the number of qubits of the chosen quantum system. Second, the projected coordinates are sequentially injected into the quantum system through simple parameterized rotations followed by a fixed, hardware-friendly entangling layer. Third, using the classical shadows strategy~\cite{huang-classical-shadows-2020}, the quantum reservoir is probed by simultaneously estimating expectation values of $k$-local observables. Finally, following a projected quantum kernel~\cite{Huang_Broughton_Mohseni_Babbush_Boixo_Neven_McClean_2021} philosophy for time series~\cite{Aaraba_Cherkaoui_Ahmad_Laprade_Nahman-Lévesque_Vieloszynski_Wang_2024}, the reservoir readout is learned within the reproducing kernel Hilbert space (RKHS) of a classical kernel function applied to the space of features generated by the quantum featurizer. 

Model capacity is explicitly characterized in a probably approximately correct (PAC) fashion~\cite{ShalevShwartz2014,mohri2018foundations}. In particular, we show that with a number of qubits lower-bounded by the data dimension (number of samples and desired time series length), one can perfectly learn a dataset constructed out of window subsequences of the studied time series. We follow this with an analysis of the generalization capacity of the model conducted on weakly dependent data. In this context, we show that the model exhibits good generalization on unseen data as we increase the number of dataset samples. This learning-theoretic analysis is further validated empirically, showing that the theory matches the empirical results. In summary, our main contributions are as follows:
\begin{enumerate}
    \item We introduce an end-to-end quantum reservoir kernel learner (\textsc{QuaRK}) that couples a hardware-realistic quantum reservoir featurizer with an RKHS kernel readout for temporal learning, while decoupling quantum resources from the data dimension to enable scalability.
    \item We adopt an RKHS-based readout (kernel ridge regression-style) that enables fast closed-form training and interpretable regularization control, applied to a feature space generated via efficient measurement of the quantum reservoir featurizer using classical shadows.
    \item We provide learning-theoretic guarantees for dependent temporal data that relate resource/design choices (projection dimension, reservoir size, multiplexing, measurement budget, and regularization) to finite-sample performance. This is further validated empirically on representative temporal learning tasks.
\end{enumerate}

%\TBDcomment{[In what follows, Section~\ref{sec:framework} ..., ]}

In what follows, Section~\ref{sec:framework} introduces the learning setup and background on temporal processes and quantum reservoirs, Section~\ref{sec:methodology} presents the \textsc{QuaRK} architecture and theoretical analysis, Section~\ref{sec:experiments} reports empirical validation on $\beta$-mixing time series tasks, and Section~\ref{sec:conslusion} discusses conclusions and future directions.
